\documentclass[11pt]{article}
\usepackage[margin=1.05in]{geometry}
\usepackage{amsmath,amssymb,amsthm,mathtools}
\usepackage{enumitem}
\usepackage[colorlinks,linkcolor=blue!55!black,citecolor=blue!55!black]{hyperref}

\makeatletter
%% protocol-header.tex -- PROTOCOL §2.3 page-1 block, as a TEMPLATE not an instance.
%%
%% PROTOCOL §2.3: "Every PDF must state, on its first page: author, date,
%% recipient, a title, and the commit hash of the work-in-progress push it
%% corresponds to.  A PDF with no commit hash cannot be checked against the
%% record and should be treated as a draft opinion, not a result."
%%
%% Written 2026-10-16 because the fault had recurred twice and I had twice
%% fixed only the instance:
%%   - 2026-10-14: the Q473 PDF went to Rick with no recipient and no hash.
%%   - 2026-10-14: the Q460 PDF went to Rick saying "Recipient: Robin Langer".
%%     He noticed, and took the commit hash as authority instead.
%% Both times I repaired the one file and wrote "fix the template, not just the
%% instance" in my notes.  This file is that fix.  \input it; do not retype it.
%%
%% USAGE, immediately after \maketitle:
%%   %% protocol-header.tex -- PROTOCOL §2.3 page-1 block, as a TEMPLATE not an instance.
%%
%% PROTOCOL §2.3: "Every PDF must state, on its first page: author, date,
%% recipient, a title, and the commit hash of the work-in-progress push it
%% corresponds to.  A PDF with no commit hash cannot be checked against the
%% record and should be treated as a draft opinion, not a result."
%%
%% Written 2026-10-16 because the fault had recurred twice and I had twice
%% fixed only the instance:
%%   - 2026-10-14: the Q473 PDF went to Rick with no recipient and no hash.
%%   - 2026-10-14: the Q460 PDF went to Rick saying "Recipient: Robin Langer".
%%     He noticed, and took the commit hash as authority instead.
%% Both times I repaired the one file and wrote "fix the template, not just the
%% instance" in my notes.  This file is that fix.  \input it; do not retype it.
%%
%% USAGE, immediately after \maketitle:
%%   %% protocol-header.tex -- PROTOCOL §2.3 page-1 block, as a TEMPLATE not an instance.
%%
%% PROTOCOL §2.3: "Every PDF must state, on its first page: author, date,
%% recipient, a title, and the commit hash of the work-in-progress push it
%% corresponds to.  A PDF with no commit hash cannot be checked against the
%% record and should be treated as a draft opinion, not a result."
%%
%% Written 2026-10-16 because the fault had recurred twice and I had twice
%% fixed only the instance:
%%   - 2026-10-14: the Q473 PDF went to Rick with no recipient and no hash.
%%   - 2026-10-14: the Q460 PDF went to Rick saying "Recipient: Robin Langer".
%%     He noticed, and took the commit hash as authority instead.
%% Both times I repaired the one file and wrote "fix the template, not just the
%% instance" in my notes.  This file is that fix.  \input it; do not retype it.
%%
%% USAGE, immediately after \maketitle:
%%   \input{protocol-header}
%%   \protocolheader{Rick (grandpa-rick)}{clio-vega/proofs}{60e4ce3}
%%
%% The recipient argument is MANDATORY and has no default on purpose: a default
%% is what produced "Recipient: Robin Langer" on a PDF addressed to Rick.

%% \maketitle CLEARS \@date (and \@author, \@title) in the article class, so a
%% macro called AFTER \maketitle that reads \@date prints an EMPTY row.  Caught
%% 2026-10-16 only because I ran pdftotext on page 1 instead of trusting rc=0:
%% the build was clean and the Date field was blank.  Hence: snapshot the date
%% HERE, at \input time in the preamble, where \@date is still live.
\let\protocol@date\@date

\newcommand{\protocolheader}[3]{%
  \begin{center}
  \fbox{\begin{minipage}{0.92\textwidth}
  \small
  \textbf{PROTOCOL §2.3 record.}\\[1mm]
  \begin{tabular}{@{}ll@{}}
  Author     & Clio \\
  Date       & \protocol@date \\
  Recipient  & #1 \\
  Repository & \texttt{#2} \\
  Commit     & \texttt{#3} \\
  \end{tabular}\\[1mm]
  This PDF corresponds to the pushed commit above.  Where this text and the
  commit disagree, \textbf{the commit is authoritative}.
  \end{minipage}}
  \end{center}
  \vspace{2mm}%
}

%%   \protocolheader{Rick (grandpa-rick)}{clio-vega/proofs}{60e4ce3}
%%
%% The recipient argument is MANDATORY and has no default on purpose: a default
%% is what produced "Recipient: Robin Langer" on a PDF addressed to Rick.

%% \maketitle CLEARS \@date (and \@author, \@title) in the article class, so a
%% macro called AFTER \maketitle that reads \@date prints an EMPTY row.  Caught
%% 2026-10-16 only because I ran pdftotext on page 1 instead of trusting rc=0:
%% the build was clean and the Date field was blank.  Hence: snapshot the date
%% HERE, at \input time in the preamble, where \@date is still live.
\let\protocol@date\@date

\newcommand{\protocolheader}[3]{%
  \begin{center}
  \fbox{\begin{minipage}{0.92\textwidth}
  \small
  \textbf{PROTOCOL §2.3 record.}\\[1mm]
  \begin{tabular}{@{}ll@{}}
  Author     & Clio \\
  Date       & \protocol@date \\
  Recipient  & #1 \\
  Repository & \texttt{#2} \\
  Commit     & \texttt{#3} \\
  \end{tabular}\\[1mm]
  This PDF corresponds to the pushed commit above.  Where this text and the
  commit disagree, \textbf{the commit is authoritative}.
  \end{minipage}}
  \end{center}
  \vspace{2mm}%
}

%%   \protocolheader{Rick (grandpa-rick)}{clio-vega/proofs}{60e4ce3}
%%
%% The recipient argument is MANDATORY and has no default on purpose: a default
%% is what produced "Recipient: Robin Langer" on a PDF addressed to Rick.

%% \maketitle CLEARS \@date (and \@author, \@title) in the article class, so a
%% macro called AFTER \maketitle that reads \@date prints an EMPTY row.  Caught
%% 2026-10-16 only because I ran pdftotext on page 1 instead of trusting rc=0:
%% the build was clean and the Date field was blank.  Hence: snapshot the date
%% HERE, at \input time in the preamble, where \@date is still live.
\let\protocol@date\@date

\newcommand{\protocolheader}[3]{%
  \begin{center}
  \fbox{\begin{minipage}{0.92\textwidth}
  \small
  \textbf{PROTOCOL §2.3 record.}\\[1mm]
  \begin{tabular}{@{}ll@{}}
  Author     & Clio \\
  Date       & \protocol@date \\
  Recipient  & #1 \\
  Repository & \texttt{#2} \\
  Commit     & \texttt{#3} \\
  \end{tabular}\\[1mm]
  This PDF corresponds to the pushed commit above.  Where this text and the
  commit disagree, \textbf{the commit is authoritative}.
  \end{minipage}}
  \end{center}
  \vspace{2mm}%
}

\makeatother

\newtheorem{thm}{Theorem}
\newtheorem{prop}[thm]{Proposition}
\newtheorem{lem}[thm]{Lemma}
\newtheorem{cor}[thm]{Corollary}
\theoremstyle{definition}
\newtheorem{defn}[thm]{Definition}
\theoremstyle{remark}
\newtheorem{rem}[thm]{Remark}

\DeclareMathOperator{\Inv}{Inv}
\DeclareMathOperator{\ord}{ord}
\newcommand{\At}{\widetilde{A}}
\newcommand{\bI}{\mathbf{I}}
\newcommand{\bM}{\mathbf{M}}
\newcommand{\sm}{\setminus}
\newcommand{\ES}{\mathcal{E}}

\title{The boundary diagonal in type $B$:\\
\normalsize $\ord_S\At^{B}_{\bI,\bI}=[\,m\in\bI\,]$, and why the order is one}
\author{Clio}
\date{2026-10-21}

\begin{document}
\maketitle
\protocolheader{Rick (grandpa-rick)}{clio-vega/proofs}{958a20a}

\begin{abstract}
I close G3 of \texttt{2026-10-19-c2-sa-divisors-full-q.tex}: the boundary-diagonal
property (BD) in type $B$, previously measured ($34/34$) and not proved.  The proof is
three lemmas and is uniform in type; the statement it proves is strictly more general
than the one asked for:
\[
  \ord_{S^{\varepsilon}_m}\At^{X}_{\bI,\bI}
  \;=\;
  [\,m\in\bI\,]\cdot\#\{\,a>0 \;:\; a\,t_m\in\Delta^{+}(X_N)\ \text{and}\ \varepsilon^{a}=1\,\}.
\]
Three consequences.  (1) \textbf{G3 is closed}, so everything in \S\S5--7 of the
10-19 paper is now proved in \emph{both} types.  (2) \textbf{G4 is derived rather than
measured}: type $B$ has no $\varepsilon=-1$ component because its boundary root is
$t_m$ and not $2t_m$, and I give the exact failure pattern.  (3) ``Order exactly one''
\emph{is} reducedness of the ambient root system --- in the non-reduced $BC_N$, where
$t_m$ and $2t_m$ are both roots, the same computation returns $2$.
I also \textbf{retract a hypothesis of my own statement}: (BD) needs no genericity
condition.  Pre-registration \textbf{P1 came back SPLIT}, and the split is sharp: the
two halves of (BD) carry \emph{complementary} hypotheses, each inert in the other half.
\end{abstract}

%=====================================================================
\section{What was open, and what the gap actually was}

The 10-19 paper fixes $m\in[N]$ and $\varepsilon$ with $\varepsilon^{h}=1$, sets
$S=S^{\varepsilon}_m:=\{w_m=\varepsilon\}$, and uses two properties of Motegi's
normalised weight functions $\At^{X}_{\bM,\bI}:=\widetilde{WX}_{N,n,r}
(\mathbf{w}_{\bM}|\mathbf{w}|\bI)$ of \texttt{arXiv:2610.15344}:
(BC), the boundary collapse, and
\[
  \textbf{(BD)}\qquad \ord_S\At^{X}_{\bI,\bI}=[\,m\in\bI\,]\quad\text{for every }\bI .
\]
(BC) is proved in both types.  (BD) is proved in type $C$ (inside \texttt{prop:short}
of \texttt{2026-10-17-c2}) and was recorded as \textbf{G3}, ``measured, not proved'',
in type $B$ --- $34/34$ at $N\le3$.  Everything in \S\S5--7 was accordingly graded
\emph{proved in type $C$, conditional on G3 in type $B$}.

\begin{rem}[the gap was a reading gap, and the locator says so]
\label{rem:locator}
The G3 note points at ``the type-$B$ diagonal is \texttt{normalizationWB} (L1386)''.
L1386 is the \emph{unexpanded} statement
$\prod_{\alpha\in\Inv(\lambda)}(1-e^{-\alpha})\prod_{\alpha\in T(\lambda)}(1-qe^{+\alpha})$
--- which is \emph{verbatim the same formula} as the type-$C$ one at L1078.  The
explicit type-$B$ expansion, with the root system written out, is at
\textbf{L1305--1335}, some eighty lines \emph{earlier}, in a subsection whose opening
sentence is ``The proofs of the type $B$ statements are essentially the same as the
corresponding type $C$ arguments\ldots For brevity, we omit these proofs'' (L1281).
I had never opened L1305--1335.  My (M4) instead carried a \emph{parameterised guess},
``the boundary term $(1-w_i^{-h})$ occurs''; the guess was right, which is exactly why
it survived.  Once the region is read, the proof below is three lemmas.
This is the same shape as the broken-locator finding of 2026-10-18: a citation that
lands next to the content rather than on it is evidence that the neighbourhood was
never read, and the repair is to read around it, not to patch the reference.
\end{rem}

%=====================================================================
\section{Setup, read first-hand from the source}

Write $h=h_X$ for the \emph{boundary coefficient}: $h_B=1$, $h_C=2$.  From L1288--1292
(type $B$) and L960--964 (type $C$),
\[
  \Delta^{+}(X_N)=\{\,t_k\pm t_j \;:\; 1\le j<k\le N\,\}\;\cup\;\{\,h\,t_j \;:\; 1\le j\le N\,\},
\]
so the two systems differ \emph{only} in the boundary root, $t_j$ versus $2t_j$.  Set
$e^{t_j}:=w_j$, extended multiplicatively, so $e^{\sum_j a_jt_j}=\prod_j w_j^{a_j}$.

For $\bI\subseteq[N]$ with associated strict partition $\lambda=\lambda(\bI)$, Motegi
computes the inversion and ascent sets explicitly (L998--1003, L1013--1019 type $C$;
L1305--1312, L1322--1328 type $B$).  In both types, \emph{with the same combinatorics
and only $h$ changing}:
\begin{align}
  \Inv(\lambda)&=\{h t_k : k\in\bI\}
   \cup\{t_j+t_k : j<k,\ j,k\in\bI\}
   \cup\{t_k-t_j : j<k,\ j\notin\bI,\ k\in\bI\},
  \label{eq:inv}\\
  T(\lambda)&=\{h t_k : k\notin\bI\}
   \cup\{t_j+t_k : j<k,\ j,k\notin\bI\}
   \cup\{t_k-t_j : j<k,\ j\in\bI,\ k\notin\bI\},
  \label{eq:T}
\end{align}
and the diagonal value is (\texttt{normalizationWC} L1078, \texttt{normalizationWB} L1386)
\begin{equation}
  \At^{X}_{\bI,\bI}
  =\prod_{\alpha\in\Inv(\lambda)}\bigl(1-e^{-\alpha}\bigr)
   \prod_{\alpha\in T(\lambda)}\bigl(1-q\,e^{+\alpha}\bigr).
  \label{eq:diag}
\end{equation}

\begin{defn}[the valuation]\label{def:ord}
Let $F=\mathbb{C}(q)(w_1,\dots,w_N)$ and $R=\mathbb{C}(q)[w^{\pm}]$.  Fix $m\in[N]$ and
$\varepsilon\in\mathbb{C}^{\times}$.  Since $R/(w_m-\varepsilon)\cong
\mathbb{C}(q)[w_j^{\pm}:j\ne m]$ is a domain (here $\varepsilon\ne0$ is used),
$(w_m-\varepsilon)$ is a prime of $R$; let $\mathcal{O}\subset F$ be the local ring,
$t:=w_m-\varepsilon$ the uniformiser, and
$k=\mathcal{O}/(t)\cong\mathbb{C}(q)(w_j:j\ne m)$ the residue field.  Write $\ord_S$ for
the associated divisorial valuation on $F^{\times}$ and $\bar x\in k$ for the residue of
$x\in\mathcal{O}$.
\end{defn}

The one structural fact used repeatedly: $\ord_S$ is a valuation, so it is
\emph{additive} on products, and $\ord_S x=0$ exactly when $x\in\mathcal{O}^{\times}$,
i.e.\ when $\bar x\ne0$ in $k$.

%=====================================================================
\section{The unit criterion}

Everything in \eqref{eq:diag} is of the form $1-c\,M$ with $M$ a Laurent monomial.  One
lemma handles all of them.

\begin{lem}[unit criterion]\label{lem:unit}
Let $M=\prod_{j}w_j^{a_j}$ be a Laurent monomial and $c\in\mathbb{C}(q)^{\times}$.  Then
$1-cM\in\mathcal{O}$, and
\begin{enumerate}[label=\textup{(\roman*)},leftmargin=2.2em]
\item $\ord_S(1-cM)\ge1$ \ $\iff$ \ $a_j=0$ for all $j\ne m$ \ and \
  $c\,\varepsilon^{a_m}=1$;
\item if moreover $a:=a_m\ne0$, then $\ord_S(1-cM)=1$ \emph{exactly}, and the cofactor
  of $t$ at the generic point of $S$ is $a\,\varepsilon^{a-1}$.
\end{enumerate}
\end{lem}

\begin{proof}
$M\in R$ and $w_j$ is a unit of $\mathcal{O}$ for every $j$ (for $j=m$ because
$\varepsilon\ne0$), so $1-cM\in\mathcal{O}$.  Its residue in $k$ is
\[
  \overline{1-cM}=1-c\,\varepsilon^{a_m}\prod_{j\ne m}w_j^{a_j}.
\]
The $w_j$ with $j\ne m$ are algebraically independent over $\mathbb{C}(q)$, so this
element of $k$ vanishes if and only if the monomial $\prod_{j\ne m}w_j^{a_j}$ is the
constant monomial and the scalar is $1$: that is, $a_j=0$ for $j\ne m$ and
$c\varepsilon^{a_m}=1$.  This is (i).

For (ii), the hypotheses of (i) give $1-cM=1-cw_m^{a}$ with $c=\varepsilon^{-a}$.
Suppose first $a>0$.  Then
\[
  1-\varepsilon^{-a}w_m^{a}
  =-\varepsilon^{-a}\bigl(w_m^{a}-\varepsilon^{a}\bigr)
  =-\varepsilon^{-a}\,(w_m-\varepsilon)\sum_{i=0}^{a-1}w_m^{i}\varepsilon^{a-1-i},
\]
and the sum has residue $a\,\varepsilon^{a-1}$, which is nonzero in $k$ because
$\varepsilon\ne0$ and $a\ne0$ in $\mathbb{C}$.  Hence $t$ divides $1-cM$ exactly once.
If $a<0$, write $1-\varepsilon^{-a}w_m^{a}=-\varepsilon^{-a}w_m^{a}
\bigl(w_m^{-a}-\varepsilon^{-a}\bigr)$ and apply the case just done to $-a>0$, the
prefactor $-\varepsilon^{-a}w_m^{a}$ being a unit of $\mathcal{O}$.
\end{proof}

\begin{rem}[where ``exactly one'' can fail, and it is not where I expected]
\label{rem:char}
The cofactor in \ref{lem:unit}(ii) is $a\varepsilon^{a-1}$, so ``order \emph{exactly}
one'' is precisely the statement that $a\ne0$ \emph{in the base field}.  Over
$\mathbb{C}(q)$ this is free.  It is worth recording that it is not free in general:
\[
  \frac{1-v^{-1}}{v-1}=\frac1v \ \xrightarrow{\,v=1\,}\ 1,
  \qquad
  \frac{1-v^{-2}}{v-1}=\frac{v+1}{v^{2}} \ \xrightarrow{\,v=1\,}\ 2 .
\]
So the type-$B$ half of this step is \emph{characteristic-free}, while the type-$C$ half
needs $\operatorname{char}\ne2$: in characteristic $2$, $1-v^{-2}=(1-v^{-1})^{2}$ and the
type-$C$ boundary factor has order $2$.  The type that was already proved is the
fragile one.
\end{rem}

%=====================================================================
\section{The two blocks}

\begin{lem}[the ascent block consists of units]\label{lem:T}
For every $\alpha\in T(\lambda)$, $\ \ord_S\bigl(1-q\,e^{\alpha}\bigr)=0$, provided
$q\ne\varepsilon^{-h}$.  If $m\notin\bI$ the hypothesis $q\ne\varepsilon^{-h}$ is needed
and used exactly once; if $m\in\bI$ it is not used at all.
\end{lem}

\begin{proof}
Apply Lemma~\ref{lem:unit}(i) with $c=q$ to each of the three families of \eqref{eq:T}.
\begin{itemize}[leftmargin=1.6em]
\item $\alpha=t_k\pm t_j$ with $j<k$: the exponent vector is supported on the two
  \emph{distinct} indices $j\ne k$, so at least one index $\ne m$ has a nonzero
  exponent.  Criterion (i) fails; the factor is a unit.  No hypothesis used.
\item $\alpha=h t_k$ with $k\notin\bI$, $k\ne m$: supported on $k\ne m$; a unit.  No
  hypothesis used.
\item $\alpha=h t_m$, which by \eqref{eq:T} occurs if and only if $m\notin\bI$:
  the residue is $1-q\varepsilon^{h}$, which is nonzero precisely when
  $q\ne\varepsilon^{-h}$.
\end{itemize}
In $F$, $q$ is transcendental over $\mathbb{C}$ and $\varepsilon^{-h}\in\mathbb{C}$, so
$q\ne\varepsilon^{-h}$ holds automatically; the point of stating it is the last bullet,
which is the \emph{only} place in the whole argument where a condition on $q$ appears,
and it appears only when $m\notin\bI$.
\end{proof}

\begin{lem}[the inversion block has at most one non-unit]\label{lem:Inv}
Let $\alpha\in\Inv(\lambda)$.  Then $\ord_S(1-e^{-\alpha})=0$ unless
$\alpha=h\,t_m$, which by \eqref{eq:inv} occurs if and only if $m\in\bI$.  In that case
$1-e^{-\alpha}=1-w_m^{-h}$ and
\[
  \ord_S\bigl(1-w_m^{-h}\bigr)=
  \begin{cases}
    1,&\varepsilon^{h}=1,\\
    0,&\varepsilon^{h}\ne1 .
  \end{cases}
\]
\end{lem}

\begin{proof}
The three families of \eqref{eq:inv} again.  For $\alpha=t_j+t_k$ ($j<k$, both in $\bI$)
and $\alpha=t_k-t_j$ ($j<k$, $j\notin\bI$, $k\in\bI$) the exponent vector of
$e^{-\alpha}$ is supported on two distinct indices, so Lemma~\ref{lem:unit}(i) fails and
the factor is a unit.  For $\alpha=h t_k$ with $k\in\bI$: $e^{-\alpha}=w_k^{-h}$ is
supported on $k$ alone, so (i) requires $k=m$ and $\varepsilon^{-h}=1$.  When $k=m$ and
$\varepsilon^{h}=1$, part (ii) applies with $a=-h\ne0$ and gives order exactly $1$;
when $\varepsilon^{h}\ne1$ the residue $1-\varepsilon^{-h}$ is a nonzero scalar.
\end{proof}

%=====================================================================
\section{The theorem}

\begin{thm}[boundary diagonal, uniform in type]\label{thm:BD}
Let $X\in\{B,C\}$, $N\ge1$, $\bI\subseteq[N]$, $m\in[N]$, and
$\varepsilon\in\mathbb{C}^{\times}$.  Then
\begin{equation}
  \ord_{S^{\varepsilon}_m}\At^{X}_{\bI,\bI}
  \;=\;[\,m\in\bI\,]\cdot
  \#\bigl\{\,a>0 \;:\; a\,t_m\in\Delta^{+}(X_N)\ \text{and}\ \varepsilon^{a}=1\,\bigr\}.
  \label{eq:master}
\end{equation}
In particular, if $\varepsilon^{h}=1$ then $\ord_S\At^{X}_{\bI,\bI}=[\,m\in\bI\,]$,
which is \textup{(BD)}; and this holds in \textbf{type $B$ for all $N$ and all $\bI$},
closing \textup{G3}.
\end{thm}

\begin{proof}
By \eqref{eq:diag} and additivity of the valuation,
\[
  \ord_S\At^{X}_{\bI,\bI}
  =\sum_{\alpha\in\Inv(\lambda)}\ord_S\bigl(1-e^{-\alpha}\bigr)
  +\sum_{\alpha\in T(\lambda)}\ord_S\bigl(1-qe^{\alpha}\bigr),
\]
every summand being defined because every factor lies in $\mathcal{O}$
(Lemma~\ref{lem:unit}).  By Lemma~\ref{lem:T} the second sum is $0$.  By
Lemma~\ref{lem:Inv} the first sum has at most one nonzero summand, contributed by
$\alpha=h t_m$, which belongs to $\Inv(\lambda)$ iff $m\in\bI$ and then contributes
$1$ iff $\varepsilon^{h}=1$.  Since $\Delta^{+}(X_N)$ contains exactly one positive
multiple of $t_m$, namely $h t_m$, the set counted in \eqref{eq:master} is $\{h\}$ if
$\varepsilon^{h}=1$ and $\varnothing$ otherwise, which is the same count.
\end{proof}

\begin{cor}[G4 is derived, not measured]\label{cor:G4}
In type $C$, $h=2$ and both $\varepsilon=1$ and $\varepsilon=-1$ satisfy
$\varepsilon^{h}=1$: the divisors $\{w_m=1\}$ and $\{w_m=-1\}$ \emph{both} carry
$\ord=[\,m\in\bI\,]$.  In type $B$, $h=1$, so $\varepsilon^{h}=1$ forces
$\varepsilon=1$; on $\{w_m=-1\}$ Theorem~\ref{thm:BD} gives
$\ord_S\At^{B}_{\bI,\bI}=0$ for \emph{every} $\bI$, so \textup{(BD)} in the form
$[\,m\in\bI\,]$ fails there, and fails on \emph{precisely} the pairs $(\bI,m)$ with
$m\in\bI$, each by exactly $1$.
\end{cor}

The 10-19 paper recorded G4 as a measured asymmetry (``$\{w_m=-1\}$ is not a root
hyperplane of $B_N$''), with the measurement taken on \emph{(BC)}.
Corollary~\ref{cor:G4} is the matching statement for \emph{(BD)}, and it is now a
consequence of one mechanism rather than two observations: $t_m$ cuts out one
hyperplane, $2t_m$ cuts out two.

\begin{cor}[``order one'' \emph{is} reducedness]\label{cor:red}
The order in \eqref{eq:master} is the number of positive roots proportional to $t_m$
that vanish on $S$.  For the reduced systems $B_N$ and $C_N$ this is at most one.  In
the non-reduced system $BC_N$, where $t_m$ \emph{and} $2t_m$ are both positive roots,
the same computation returns $\ord_S\At_{\bI,\bI}=2$ at $\varepsilon=1$ and $1$ at
$\varepsilon=-1$, for every $\bI\ni m$.
\end{cor}

\begin{rem}[why this matters downstream]
(BD) enters the 10-19 paper through
$\ord_S\det\Phi=\sum_{\bI}\ord_S\At_{\bI,\bI}=\#\{\bI\ni m\}=2^{N-1}$ and through
$\ord_S(\Phi^{-1})_{\nu\bM}=-1$.  Both are \emph{colength} statements, so
Corollary~\ref{cor:red} says the boundary colength of my lattice is a reducedness
count.  A non-reduced ambient system would double it.  I flag this because the live
question Q538 asks whether my object's surviving parity is $\bmod 4$ rather than
$\bmod 2$ on the grounds that the relevant affine type $A^{(2)}_{2\ell}$ is the
\emph{non-reduced} one --- and \eqref{eq:master} is a place where non-reducedness
produces an off-by-a-factor-of-two in a quantity I actually use.  I am not claiming the
two are the same phenomenon; I am recording that the discriminator is the same
($\alpha$ and $2\alpha$ both roots) and that here it is now computed, not conjectured.
\end{rem}

%=====================================================================
\section{A hypothesis of my own statement, retracted}

\begin{rem}[(BD) needs no genericity: \texttt{rem:generic} is withdrawn]
\label{rem:retract}
The 10-19 statement of (BD) reads ``$\ord_S\At^{X}_{\bI,\bI}=[\,m\in\bI\,]$ for every
$\bI$, \emph{at a point of $S$ generic in the sense of Remark \texttt{rem:generic}}'', and
that remark explains: $\At_{\bI,\bI}$ contains factors $(1-qw_mw_k)$ and $(1-qw_k/w_m)$
which at $w_m=\varepsilon$ become $(1-q\varepsilon w_k)$ and
$(1-q\varepsilon^{-1}w_k)$, vanishing on the hypersurfaces $w_k=\pm q^{-1}$; so
``generic'' was defined to mean \emph{outside} those hypersurfaces.

\textbf{That hedge is not a hypothesis of the theorem, and I withdraw it.}  It records a
limitation of how I was \emph{measuring}, not a condition on what is \emph{true}.
$\ord_S$ is the valuation at the generic point of the divisor $S$
(Definition~\ref{def:ord}): the residue field is $k=\mathbb{C}(q)(w_j:j\ne m)$, in which
$1-q\varepsilon w_k$ is a nonzero element because $w_k$ is an indeterminate.  The
hypersurfaces $\{w_k=\pm q^{-1}\}$ have codimension $2$ in the ambient space and
codimension $1$ \emph{inside} $S$; they are proper closed subsets of $S$ and are
therefore invisible to a divisorial valuation.  They are visible only if one evaluates
at a \emph{numerical} point of $S$ and calls the result $\ord_S$ --- which is what my
first $N=3$ census did when it drew $w_3=7/3$ with $q=3/7$ and reported
$\ord_S\det\Phi=8$ against the correct $4$.

The honest description of that episode is therefore not ``(BD) requires a generic
point'' but ``my instrument required a generic point, and I wrote the instrument's
limitation into the theorem''.  Every computation in the present paper keeps $q$ and
$w_{j\ne m}$ \emph{symbolic} and so sits at the generic point by construction; the
hypersurfaces cannot be hit, and no sampling is needed.  The replacement hypothesis is
the sharp one supplied by Lemma~\ref{lem:T}: $q\ne\varepsilon^{-h}$, and only on the
$m\notin\bI$ half.
\end{rem}

%=====================================================================
\section{Pre-registration: which arm fired}
\label{sec:prereg}

Written before computing, in
\texttt{scratch/20261021c2/PREREG.md}; graded arm by arm.

\begin{description}[leftmargin=1.4em]
\item[P1 --- ``(BD) in type $B$ holds for all $N$, all $\bI$, by an argument uniform in
  type''.]  \textbf{SPLIT}, as predicted, and the split is a result rather than a
  blemish.  The theorem is true and the argument is uniform; what is \emph{not} uniform
  is the hypothesis class, and the two halves of (BD) turn out to carry
  \textbf{complementary} hypotheses:
  \begin{center}
  \begin{tabular}{lll}
   \textbf{half} & \textbf{needs} & \textbf{does not need}\\\hline
   $m\in\bI$ \ (value $1$) & $h\ne0$ in the base field & any condition on $q$\\
   $m\notin\bI$ \ (value $0$) & $q\ne\varepsilon^{-h}$ & any condition on $h$\\
  \end{tabular}
  \end{center}
  Neither hypothesis is visible from the other half.  The $m\in\bI$ half is controlled
  by $\Inv(\lambda)$ (Lemma~\ref{lem:Inv}), the $m\notin\bI$ half by the single
  boundary element of $T(\lambda)$ (Lemma~\ref{lem:T}), and $h t_m$ lies in exactly one
  of the two sets.  This is the \emph{second} time in two sessions that a lemma's two
  halves separated this way (2026-10-21 c1: $k<t$ irremovable in one half of the
  accounting bijection, inert in the other), and the \emph{third} counting the S/A
  grouping of 10-19.  A binary register cannot record any of them.
\item[P2 --- ``no genericity hypothesis is needed''.]  \textbf{TRUE}; see
  Remark~\ref{rem:retract}.  This is a retraction of a hypothesis I myself attached to
  (BD), so it had no natural objector.
\item[P3 --- ``the order is exactly $1$ because $B_N$ and $C_N$ are reduced''.]
  \textbf{TRUE}, and the pre-registered falsifier was run: in the non-reduced $BC_N$ the
  order is $2$ at $\varepsilon=1$ (Corollary~\ref{cor:red}), confirmed
  computationally.  The same mechanism delivers G4 (Corollary~\ref{cor:G4}).
\end{description}

%=====================================================================
\section{Verification}
\label{sec:verif}

All computations are pure Python $+$ SymPy (\textbf{SageMath is not installed}), with
$q$ and $w_{j\ne m}$ kept \emph{symbolic}: $\ord_S$ is computed by substituting
$w_m\mapsto\varepsilon+t$ and taking the $t$-adic valuation of the resulting rational
function, which is the valuation at the generic point of $S$.  Scripts:
\texttt{scratch/20261021c2/\{g3.py,run\_main.py,run2.py,run3.py,run4.py\}}.

\begin{description}[leftmargin=1.4em]
\item[Control 1 --- the transcription.]  My transcription of the explicit type-$B$
  Inv/$T$ products (L1305--1335) agrees with $\At^{X}_{\bI,\bI}$ computed from Motegi's
  \emph{$S_n$ symmetrisation definition} --- an independent mechanism which never
  mentions a positive root --- for \textbf{$60/60$} pairs $(\bI,X)$ at $N\le4$, both
  types.  The engine is \texttt{scratch/q484\_engine.py} carrying the 2026-10-17
  $\mathcal{E}^{B}$ correction (before that date it divided $WB$ by $\mathcal{E}^{C}$).
\item[Control 2 --- the factor list.]  The list of factors indexed by
  $\Inv(\lambda)\sqcup T(\lambda)$ multiplies back to the diagonal, $60/60$ at $N\le4$,
  both types.  So no factor is dropped or double-counted in the bookkeeping that
  Lemmas~\ref{lem:T} and \ref{lem:Inv} partition.
\item[Control 3 --- the instrument is sensitive.]  Four deliberate mis-transcriptions
  were run against the engine and \emph{all four were detected}: \texttt{h\_swap}
  (type-$C$ boundary root in type $B$) $10/10$; \texttt{pair\_leq} (non-strict Inv pair
  product) $10/10$; \texttt{drop\_bdry} $10/10$; \texttt{T\_h\_swap} $8/10$, the two
  misses being $\bI=[N]$, where $\bI^{c}=\varnothing$ and the plant is a no-op --- i.e.\
  correct behaviour, not a failure.
\item[Measurement A --- the theorem.]  $\ord_S\At^{X}_{\bI,\bI}$ measured from the
  symmetrisation value (so using neither Lemma~\ref{lem:unit} nor the factorisation),
  for all $\bI$ and all $m$:
  \begin{center}
  \begin{tabular}{llll}
   & $N=1$ & $N=2$ & $N=3$\\\hline
   type $B$, $\varepsilon=+1$ & $2/2$ & $8/8$ & $24/24$\\
   type $C$, $\varepsilon=+1$ & $2/2$ & $8/8$ & $24/24$\\
   type $C$, $\varepsilon=-1$ & $2/2$ & $8/8$ & $24/24$\\
   type $B$, $\varepsilon=-1$ & $1/2$ & $4/8$ & $12/24$\\
  \end{tabular}
  \end{center}
  The type-$B$, $\varepsilon=-1$ row is Corollary~\ref{cor:G4} and its disagreements are
  \emph{exactly} the $N\cdot2^{N-1}=1,4,12$ pairs with $m\in\bI$, every one measured $0$
  against a predicted $1$.  The $34/34$ of the top row reproduces the 10-19 G3 number on
  the same domain.
\item[Measurement B --- the general formula \eqref{eq:master}.]  $306/306$ over nine
  arms, $X\in\{B,C,BC\}\times\varepsilon\in\{1,-1,i\}$, at $N\le3$.  The
  $\varepsilon=i$ arm is a genuine control: there the prediction is $0$ everywhere
  (since $i^{1},i^{2}\ne1$), so the formula is not merely fitted to $\varepsilon=\pm1$.
\item[Measurement C --- the reducedness lesion.]  Non-reduced $BC_N$, $N\le4$:
  $\ord=2$ on all $49$ cases with $m\in\bI$ at $\varepsilon=1$, and $\ord=1$ on all $49$
  at $\varepsilon=-1$; $\ord=0$ on all $49$ cases with $m\notin\bI$ in both.
\item[Measurement D --- the split lesion.]  Setting $q:=1=\varepsilon^{-h}$ at
  $\varepsilon=1$, in violation of Lemma~\ref{lem:T}'s hypothesis: the $m\in\bI$ half is
  \textbf{inert} ($1\to1$, $129/129$) and the $m\notin\bI$ half \textbf{breaks}
  ($0\to1$, $129/129$), both types, $N\le5$.  Corroborated by a second route
  (full-product expansion rather than factorwise sum), $49/49$ each at $N\le4$.  So the
  hypothesis is irremovable from one half and inert in the other, which is P1's split
  exhibited as a lesion rather than argued.
\item[Positive control, named.]  The pre-registration named $N=3$, $m=2$,
  $\bI=\{2,3\}$, type $B$, $\varepsilon=1$ and asserted that all blocks are populated.
  \textbf{That claim was wrong}: at $N=3$ with $\bI^{c}=\{1\}$ the $T$ pair family
  (needs $|\bI^{c}|\ge2$) and the $T$ difference family (needs some $j\notin\bI$ above
  an element of $\bI$) are both \emph{empty}, so only $4$ of the $6$ root families
  occur.  The case still discriminates --- it is non-vacuous ($r=2\ge1$, so not the
  $r=0$ vacuum, where a green carries zero bits) --- but it is weaker than advertised.
  The smallest case populating \emph{all six} families is $N=4$, $m=2$,
  $\bI=\{2,3\}$, $\bI^{c}=\{1,4\}$; there the diagonal has $10$ factors, exactly one of
  which ($1-w_2^{-1}$, from $\alpha=t_2\in\Inv(\lambda)$) has positive order, and
  $\ord_S=1$ from the symmetrisation agrees.  The companion $m=1\notin\bI$ gives $0$.
  Both cases are printed factor-by-factor in \texttt{run3.py}/\texttt{run4.py} output.
\end{description}

\begin{rem}[a fault in my own instrument plumbing, recorded]
The trajectory emitter's \texttt{log\_dir} defaults to the \emph{relative} path
\texttt{state/trajectory}, so this session's events first landed under
\texttt{projects/state/trajectory/} with a bare-timestamp session tag --- the hazard
documented in \texttt{emit.py}'s own docstring, which I had read and not acted on.
Consolidated to \texttt{clio-prove-G3-20261021c2.jsonl}; $18/18$ emitted events
recorded (a \texttt{verify} lacking \texttt{verification\_target} is silently dropped,
so the count is checked, not assumed).  \texttt{tracecheck} also flagged one of my
\texttt{verify} events circular, \textbf{correctly}: I had offered the
transcription-versus-engine comparison as evidence \emph{for the transcription}, and
that comparison descends from the transcription.  The corrected encoding makes the
engine the evidence and the comparison the verify event, and demotes the plant census
to what it actually certifies --- the \emph{instrument}, not the transcription, since
the plants are themselves built from it.  Fourth instance of this class; the as-recorded
file is kept as \texttt{*.asrecorded-circular-encoding}.
\end{rem}

%=====================================================================
\section{What this closes, and what remains}

\begin{description}[leftmargin=1.4em]
\item[G3 --- CLOSED.]  (BD) in type $B$ is Theorem~\ref{thm:BD}.  Consequently
  \S\S5--7 of \texttt{2026-10-19-c2-sa-divisors-full-q.tex} --- Proposition
  \texttt{prop:lat}, Lemma \texttt{lem:inverse}, Lemma \texttt{lem:resfun},
  Theorem~\texttt{thm:S} and Theorem~\texttt{thm:Sprime} --- are \textbf{proved in both
  types}, no longer ``proved in type $C$ and conditional on G3 in type $B$''.  Note the
  conditionality was on (BD) alone: (BC) in type $B$ is the gauge theorem of
  \texttt{2026-10-19-typeB} and was never in doubt.
\item[G4 --- upgraded from measured to derived.]  Corollary~\ref{cor:G4}, with the exact
  failure pattern.  The 10-19 measurement was of (BC); this is the (BD) half.
\item[G1, G5 --- already closed] (2026-10-20 c1 and 2026-10-21 c1 respectively).
\item[G2 --- still open.]  Non-cancellation at fixed $\chi$; untouched here, and
  nothing above bears on it.
\item[A new hypothesis, replacing a withdrawn one.]  $q\ne\varepsilon^{-h}$, needed only
  on the $m\notin\bI$ half.  Automatic in $F$, but it should be carried in any
  specialisation of $q$, and in particular \textbf{a $q\to1$ limit destroys (BD)'s
  dependence on $\bI$ entirely} (Measurement D): the order becomes $1$ for every $m$ and
  every $\bI$, so the colength count $\#\{\bI\ni m\}=2^{N-1}$ becomes $2^{N}$.  Any
  future classical ($q\to1$) degeneration of this material must not inherit (BD).
\end{description}

\paragraph{One thing I did not do.}
I have not re-read Motegi's \emph{proof} of \eqref{eq:diag} in type $B$; the paper omits
it (``essentially the same as the corresponding type $C$ arguments\ldots we omit these
proofs'', L1281) and derives it at L3797 from \texttt{relationwithAnormalizedtypeB} plus
\texttt{RTVtheorem}(ii).  What I have done is take \eqref{eq:diag} as the statement to
be used --- it is a theorem of his paper, read first-hand at L1386 and with its explicit
expansion read first-hand at L1305--1335 --- and verify that expansion against his
\emph{definition} by an independent mechanism ($60/60$, three controls, four plants).
So the dependency here is Motegi's type-$B$ diagonal \emph{theorem}, corroborated
against his \emph{definition}, and \textbf{not} re-derived from the RTV stable-envelope
input.  That input is the same unverified root recorded as G5 of the 10-19 paper and now
tracked as a registry node; this paper does not reduce it, and I am registering
\eqref{eq:diag} as a premise node with its own grade rather than letting a sentence in
this paragraph carry it.

\end{document}
